% !TeX root = ../main.tex
\section{Introduction}
\label{sec:introduction}

Compute-in-memory (CIM) architectures execute dense matrix-vector multiplications inside memory arrays, reducing the data movement that dominates conventional digital inference \citep{horowitz2014computing,peng2020dnnneurosim,lanza2025memristor}. The same physical proximity that makes CIM attractive also exposes neural networks to device-level constraints: finite conductance precision, device-to-device variability, cycle-to-cycle read noise, and retention drift \citep{rasch2021aihwkit,crosssim2024,xu2025emerging,guo2024organic,photonics2025organicreview}. These effects are often measured or modeled at the device level, but their system-level consequences for modern vision backbones remain difficult to rank before fabrication.

The central difficulty is diagnostic. Different failure modes can produce the same accuracy loss: a low-precision model may overfit one sampled mismatch realization during hardware-aware training, lose information to a coarse conductance grid, or drift after programming. Mature CIM simulators expose increasingly detailed non-idealities \citep{peng2020dnnneurosim,lammie2022memtorch,rasch2021aihwkit,crosssim2024}, yet device-specific studies, including organic and optoelectronic demonstrations, often treat variability, ADC resolution, retention-induced drift, and transfer across fabricated instances qualitatively or in isolation \citep{zeng2023organicmemristor,jung2024organicfilaments,lin2016physical}. A deployment workflow therefore needs to distinguish an algorithmic robustness problem from a physical precision-retention limit.

Here we present a behavioral simulation framework that evaluates analog neural network deployment under both idealized pure-quantization and realistic device-physics regimes. We apply a mixed-signal deployment model to the Tiny-ViT-5M backbone and use the contrast between these regimes to separate cross-instance generalization from physical retention limits. The paper makes three contributions. First, it maps a 4/6/8-bit precision-retention frontier for PCM and finds a non-monotonic precision--accuracy curve. In this tested sweep, 8-bit PCM is drift-flat with the highest fresh accuracy (77.60\%), 4-bit PCM is trainable but drift-limited, and 6-bit PCM enters a D2D-sensitive transition zone with lower mean accuracy (68.44\%, seed-level std 6.03\%) despite drift stability. This behavior is consistent with a quantization-D2D interaction in which mismatch perturbations become comparable to the 1/64 conductance step. Second, the study explains why this frontier cannot be predicted from algorithm-only tests: in an idealized pure-quantization regime, standard fixed-instance hardware-aware training at 4-bit collapses to near-chance, while Ensemble Hardware-Aware Training (HAT) restores fresh-instance accuracy to 86.16$\pm$0.19\% by resampling D2D mismatch masks each epoch. Third, the comparison yields a deployment rule: solve cross-instance robustness first, then choose the memory precision and readout regime that satisfy retention and variability constraints.

The framework is intentionally behavioral rather than circuit-accurate. It focuses on macroscopic precision, mismatch, readout, and drift phenomena, while leaving array-level parasitics such as IR drop and sneak paths to future circuit-calibrated studies. Additional supporting analyses, including front-end compensation and zero-shot transfer to alternative organic optoelectronic profiles, are provided in the Supplementary Information. The goal is not to replace detailed circuit simulators, but to provide a reproducible layer for comparing mitigation strategies and identifying deployment frontiers before hardware-specific calibration is available.

\begin{figure}[htbp]
\centering
\makebox[\textwidth][c]{\includegraphics[width=1.1\textwidth]{fig1_algorithmic_rescue}}
\caption{\textbf{Algorithmic generalization failure in the ideal-device ablation.} (A) Fresh-instance accuracy for the AIHWKit IdealDevice baseline. 4-bit pure quantization falls to near chance, while Ensemble HAT restores performance near the 8-bit baseline. This is an algorithmic ablation, not the PCM deployment headline. (B) Loss-landscape interpolation from the source training mask ($\alpha=0$) to fresh hardware instances ($\alpha=1$). Ensemble HAT discovers a broad plateau robust to D2D mismatch, while Standard HAT exhibits extreme sensitivity.}
\label{fig:algorithmic_rescue}
\end{figure}
